\section{Methodology}

Our investigation tests the gradient competition hypothesis through a two-stage experimental design: first establishing that spurious feature dominance exists (Existence, H-E1), then testing whether this dominance arises from stronger gradient signals (Mechanism, H-M1). This sequential design ensures we do not test a mechanism for a phenomenon that may not exist.

\subsection{Overview}

The gradient competition hypothesis posits that spurious features dominate because they produce stronger gradient signals during training. If true, we should observe:

\begin{enumerate}
    \item \textbf{Existence:} GradCAM attribution on spurious regions exceeds attribution on core regions early in training
    \item \textbf{Mechanism:} Gradient norms from spurious-aligned samples exceed those from minority samples, especially in early epochs
\end{enumerate}

We design experiments to test both predictions. Existence confirmation (H-E1) enables mechanism testing (H-M1). Mechanism falsification does not invalidate the existence of spurious dominance---it invalidates a specific explanation for why dominance occurs.

\subsection{Experiment 1: Spurious Feature Attribution (H-E1)}

\textbf{Objective:} Measure whether spurious features dominate core features early in training.

We use GradCAM \citep{Selvaraju2017GradCAM} to compute class-discriminative attributions at layer4 (final convolutional block) of ResNet-50. For each validation sample, we obtain a spatial attribution map indicating which image regions contribute to the classification decision.

We define spurious and core regions based on Waterbirds image structure:
\begin{itemize}
    \item \textbf{Spurious region:} Upper 60\% of image (background habitat)
    \item \textbf{Core region:} Lower 40\% of image (bird location)
\end{itemize}

For each epoch $t$, we compute the attribution ratio:
\begin{equation}
R(t) = \frac{\sum_{i} A_{\text{spurious}}^{(i)}(t)}{\sum_{i} A_{\text{core}}^{(i)}(t)}
\end{equation}

\textbf{Success criterion:} $R(t) > 1.0$ for $t < 10$ (spurious dominance before epoch 10).

\subsection{Experiment 2: Gradient Norm Analysis (H-M1)}

\textbf{Objective:} Test whether spurious features receive stronger gradient signals.

We partition training samples into two groups based on spurious correlation alignment:
\begin{itemize}
    \item \textbf{Spurious-aligned:} Samples where label matches background type (e.g., waterbird on water)
    \item \textbf{Minority:} Samples where label mismatches background (e.g., waterbird on land)
\end{itemize}

For each training batch, we compute gradient norms at layer4:
\begin{equation}
G_{\text{spurious}}(t) = \frac{1}{|S_t|} \sum_{i \in S_t} \|\nabla_\theta \mathcal{L}(x_i, y_i)\|_2
\end{equation}
\begin{equation}
G_{\text{minority}}(t) = \frac{1}{|M_t|} \sum_{i \in M_t} \|\nabla_\theta \mathcal{L}(x_i, y_i)\|_2
\end{equation}

We compute the gradient norm ratio: $\rho(t) = G_{\text{spurious}}(t) / G_{\text{minority}}(t)$

\textbf{Hypothesis:} If gradient competition drives spurious dominance, then $\rho(t) > 1.5$ in early epochs (1--10).

\textbf{Falsification:} If $\rho(t) < 1.0$, the mechanism is inverted---minority samples produce stronger gradients.

\subsection{Training Configuration}

Both experiments use identical training setup: Waterbirds v1.0 dataset, ResNet-50 (ImageNet pretrained), SGD optimizer (momentum=0.9, weight\_decay=1e-4), batch size 128, 50 epochs. H-E1 uses LR=0.01; H-M1 uses LR=0.001 for stable gradient measurements.

\subsection{Gate Structure}

\begin{table}[h]
\centering
\begin{tabular}{@{}llll@{}}
\toprule
Hypothesis & Gate Type & Condition & Action if Fail \\
\midrule
H-E1 & MUST\_WORK & $R(t) > 1.0$ for $t < 10$ & Abort \\
H-M1 & MUST\_WORK & $\rho(t) > 1.5$ for $t \in [1,10]$ & Pivot \\
\bottomrule
\end{tabular}
\caption{Gate structure for hypothesis testing.}
\label{tab:gates}
\end{table}
